\section{Distribución de Poisson}
\subsection{Definición}
La distribución de probabilidad de variable aleatoria $x$ es:
\begin{equation}
    P(\chi=x) = \frac{e^{-\lambda}\lambda^x}{x!} \qquad x\in \mathbb{N}
\end{equation}
donde $\lambda$ es el número promedio de resultados que ocurren en el intervalo dado.
Para esta definición se tiene que:
\begin{itemize}
    \item $\mu=\lambda$
    \item $\sigma^2= \lambda$
\end{itemize}
\subsection{Propiedades}
\begin{itemize}
    \item El número de ocurrencias en cualquier intervalo es independiente del número que se tiene en otros intervalos.
    \item La probabilidad de más de una ocurrencia en cualquier punto único es aproximadamente cero.
    \item La probabilidad de una ocurrencia es la misma a través de todo el campo de observación.
\end{itemize}
\begin{ejemplo}
    Un departamento de reparación de motores recibe un promedio de cinco llamadas de servicio por hora,¿ cual es la probabilidad que reciban tres llamadas en una hora?
    \begin{align*}
        P(X=3) & = \frac{e^{-5}5^3}{3!} \\
               & = 0.14037
    \end{align*}
\end{ejemplo}
\begin{ejemplo}
    Si un banco recibe un promedio de 6 cheques falsos al día: ¿cuáles son las probabilidades de que reciba
    \begin{itemize}
        \item cuatro cheques falsos en un día cualquiera?
              \begin{align*}
                  P(X=4) & = \frac{e^{-6}6^4}{4!} \\
                         & = 0.13385
              \end{align*}
        \item 10 cheques falsos en dos días consecutivos cualquiera?
              \begin{align*}
                  P(X=3) & = \frac{e^{-12}12^10}{10!} \\
                         & = 0.10483
              \end{align*}
    \end{itemize}
\end{ejemplo}
\begin{ejemplo}
    Supongamos que los defectos de la hilaza se pueden aproximar mediante unproceso Poisson con una media de 0.2 defectos/metro. Si se inspeccionan longitudes de 6 metros. Calcula la probabilidad de hallar menos de dos defectos.
    \begin{align*}
        P(x<2) & =P(x\leq 1)                                        \\
               & = P(x=1)+P(x=0)                                    \\
               & =\frac{e^{-1.2}1.2^1}{1!}+\frac{e^{-1.2}1.2^0}{0!} \\
               & =0.3614+0.3011                                     \\
               & = 0.6625
    \end{align*}
\end{ejemplo}
\subsection{Aproximación de la distribución de Poisson a la distribución binomial}
En determinadas circunstancias la distribución de Poisson se puede utilizar para aproximar probabilidades binomiales. La aproximación es más adecuada cuando la probabilidad del éxito p está cerca de 0. Para utilizar la aproximación, es necesario determinar solamente la media o el valor esperado de la distribución binomial. Este valor se utiliza como la media del proceso para la distribución de Poisson. Es decir, el promedio del proceso  $\mu$ es igual al promedio binomial np.
\begin{ejemplo}
    Se sabe que el 5\% de los libros encuadernados en cierto taller tienen encuadernación defectuosa. Calcula la probabilidad de que 2 de 100 libros terminados en este taller tengan una encuadernación defectuosa, utilizando:
    \begin{itemize}
        \item la fórmula para la distribución binomial.
              \begin{align*}
                  P(x) & = C\left(\begin{matrix}
                          n \\
                          x
                      \end{matrix}\right) p^x (1-p)^{n-x}      \\
                       & = C\left(\begin{matrix}
                          100 \\
                          2
                      \end{matrix}\right) (0.05)^2 (0.95)^{98} \\
                       & = 0.08118
              \end{align*}
        \item la aproximación de Poisson a la distribución binomial.
              \begin{align*}
                  P(x=2) & = \frac{e^{-5}5^2}{2!} \\
                         & = 0.0842
              \end{align*}
    \end{itemize}
\end{ejemplo}